\subsection*{一步一步建立直观}

\paragraph{1. 分别观察边界高度和已积面积}

右端点处的曲线高度只描述“此刻增加得多快”，而此前曲线下的面积描述“到此刻已累积多少”。移动终点时，哪怕已积面积很大，新增长条的高度仍只由边界附近决定。导数关注局部增量，因此不会把整个历史面积直接拿来当斜率。



\paragraph{2. 粗糙估计也能先说明结构}

减少矩形数后，估计会有明显误差，但每一块仍是高度乘宽度，物理单位也保持正确。若高度单位是米每秒、宽度是秒，面积单位就是米。单位检查让我们在正式求极限前，就能发现把速率与位移混在一起的错误。



\paragraph{3. 用误差夹逼支持直观}

对递增连续函数，左端点和与右端点和分别从下方和上方夹住面积。分割细化使二者之差趋零，直观上的“贴紧曲线”就有了定量表达。本图只画左端点和；你可以用端点高度差与区间宽度自行估计另一侧的误差界。
